\begin{table}[t]
\centering
\caption{FLC-O frente a cada referencia en la Tabla VIII. Mejora media del retardo por par de caudales (\%); IC de 95\,\% de esa media (t de Student sobre las 16 mejoras por par); rango (m\'inimo a m\'aximo) de las 16 mejoras, que no es un intervalo de confianza; y n\'umero de pares (de 16) en que FLC-O es significativamente mejor o peor (Wilcoxon pareado sobre 20 r\'eplicas, correcci\'on de Holm, 0.05).}
\label{tab:comparaciones}
\scriptsize
\setlength{\tabcolsep}{2.6pt}
\begin{tabular}{@{}lrrrrr@{}}
\toprule
Referencia & Media & IC 95\,\% & Rango & Mejor & Peor \\
\midrule
Pappis y Mamdani reconstruido & +28.2 & [+24.1, +32.4] & +19.5 a +49.3 & 16 & 0 \\
Ciclo fijo (Webster) & +44.1 & [+41.0, +47.1] & +39.2 a +61.2 & 16 & 0 \\
Actuado (hueco 3 s) & +31.3 & [+23.5, +39.1] & +11.4 a +56.8 & 16 & 0 \\
Actuado afinado & +22.0 & [+14.9, +29.1] & +2.9 a +46.5 & 16 & 0 \\
Actuado afinado con vista adelante & +7.7 & [+3.6, +11.8] & +0.5 a +26.0 & 14 & 0 \\
Óptimo (PD, caudales conocidos) & +19.2 & [+12.9, +25.5] & +2.7 a +43.5 & 16 & 0 \\
\bottomrule
\end{tabular}
\fuente{simulaci\'on del autor; Wilcoxon exacto o con aproximaci\'on normal seg\'un empates (estadistica.py).}
\end{table}
